\section{Closed-form benchmark solutions}\label{app:closedforms}

This appendix collects the closed forms that serve as acceptance targets for the
code-verification benchmarks CV-1--CV-4. Each solution is small-strain, linear-elastic,
and quasi-static. Each is also the limit a finite-deformation contact solve must reproduce
once the loading and material parameters enter the appropriate regime, so a numerical
result that misses these targets signals a coding error rather than a modeling choice. We
derive each solution symbolically in \code{docs/hertz\_derivation/} (SymPy, with printed
self-checks), and the verification harness compares against the vectorized numpy evaluators
in \code{postprocessing/contact\_fields.py}. Throughout the paper we use one sign
convention: a normal gap $\gn<0$ denotes penetration, and compressive stress is negative.

The three-dimensional axisymmetric forms below are the canonical reference, and the
numerical study restricts itself to their two-dimensional reductions. The verification in
Section~\ref{sec:verif} exercises the line-contact limit of CV-1, namely the contact
half-width, and the planar reduction of CV-2, whose stick-radius exponent is $1/2$ rather
than the axisymmetric $1/3$. The plane-strain solver reaches these regimes; we leave full
three-dimensional Hertz to future work.

\subsection{CV-1 --- Hertz normal contact}\label{app:hertz}

Two elastic bodies of radii $R_1,R_2$ and moduli $(E_1,\nu_1),(E_2,\nu_2)$ are pressed
together by a normal force $F$. Hertz theory replaces this two-body problem with a single
contact between an effective elastic medium and an effective curved surface. We therefore
collapse the four elastic constants and two radii into one modulus and one radius:
% code: postprocessing/contact_fields.py::hertz_3d_params
\begin{equation}\label{eq:hertz-effective}
  \frac{1}{E^\ast}=\frac{1-\nu_1^2}{E_1}+\frac{1-\nu_2^2}{E_2},
  \qquad
  \frac{1}{R}=\frac{1}{R_1}+\frac{1}{R_2},
\end{equation}
where $E^\ast$ is the effective contact modulus and $R$ the effective radius of the gap.
With $R$ fixed, the undeformed gap near the apex is the paraboloid $\gn(r)=r^2/2R$, with
$r$ the in-plane distance from the contact axis. The axisymmetric Hertz solution
\citep{johnson1985contact} then gives the contact radius $a$, the pressure $p(r)$, its
peak value $p_0$, the approach $\delta$ (the relative rigid-body motion of the two bodies),
and the force--approach law:
% code: postprocessing/contact_fields.py::hertz_3d_params
% code: postprocessing/contact_fields.py::hertz_pressure
\begin{align}
  a &= \left(\frac{3FR}{4E^\ast}\right)^{1/3}, \label{eq:hertz-a}\\
  p(r) &= p_0\sqrt{1-r^2/a^2}, \qquad
  p_0 = \frac{2E^\ast a}{\pi R} = \frac{3F}{2\pi a^2}, \label{eq:hertz-p}\\
  \delta &= \frac{a^2}{R}, \qquad
  F = \frac{4}{3}\,E^\ast\sqrt{R}\,\delta^{3/2}. \label{eq:hertz-force}
\end{align}
The contact radius grows as the cube root of the load, and the pressure traces a
hemispherical cap whose peak $p_0$ is $3/2$ the mean pressure $F/\pi a^2$. Integrating
this pressure over the contact recovers the applied load, $\int_0^a p(r)\,2\pi
r\,\mathrm{d}r = F$. The harness uses this identity as an independent consistency check on
the numerical pressure field.

Beneath the contact, the load redistributes into the bulk, and the stresses on the
symmetry axis follow \citep[Eq.~3.45]{johnson1985contact}:
% code: postprocessing/contact_fields.py::hertz_subsurface_axis
\begin{equation}\label{eq:hertz-subsurface}
  \sigma_z = -\frac{p_0}{1+\zeta^2}, \qquad
  \sigma_r = \sigma_\theta
  = -p_0\!\left[(1+\nu)\!\left(1-\zeta\arctan\tfrac{1}{\zeta}\right)-\tfrac12\,\frac{1}{1+\zeta^2}\right],
  \qquad \zeta=\frac{z}{a}.
\end{equation}
Here $\sigma_z$ is the axial stress, $\sigma_r=\sigma_\theta$ the in-plane normal
stress, and $\zeta=z/a$ the depth measured in contact radii. The difference between these
components sets the shear that drives yield, which is why the maximum shear sits below the
surface rather than at it. For $\nu=0.3$ the maximum shear
$\tfrac12|\sigma_z-\sigma_r|$ reaches $\tau_{\max}=0.31\,p_0$ at depth $z/a=0.481$, the
von Mises stress peaks at $\sigma_{vM}^{\max}=0.62\,p_0$, and first yield occurs at
$p_0\approx 1.61\,\sigma_Y$. We accept the numerical solution when $a$, $p_0$, $\delta$,
and $F(\delta)$ agree to $<1\%$, the pressure integral to $<0.5\%$, and the subsurface
maximum shear lies at $z/a=0.48\pm0.02$ with magnitude $0.31\,p_0\pm2\%$.

The plane-strain solver reaches the line-contact limit rather than the full axisymmetric
case, so we use that limit as the indenter benchmark for the chart-FEM verification. For a
load per unit thickness $P'$, the half-width and the peak pressure are
% code: postprocessing/contact_fields.py::line_contact_params
\begin{equation}\label{eq:line-contact}
  a = 2\sqrt{\frac{P'R}{\pi E^\ast}}, \qquad
  p_0 = \frac{2P'}{\pi a} = \sqrt{\frac{P'E^\ast}{\pi R}},
\end{equation}
with $E^\ast = E/2(1-\nu^2)$ for two identical bodies and $a$ now the contact half-width
of the strip. The half-width grows as the square root of the load here, in contrast to the
cube-root growth of the axisymmetric radius in \eqref{eq:hertz-a}. The subsurface field of
a Hertzian half-plane pressure over $[-a,a]$ is the McEwen solution
\citep[Eq.~4.49]{johnson1985contact}, implemented in \code{line\_contact\_subsurface} and
evaluated for the stress-bulb figures.

\subsection{CV-2 --- Cattaneo--Mindlin frictional partial slip}\label{app:cattaneo}

A tangential load $Q<\mu P$ is applied to a pre-formed Hertz contact of radius $a$ and
peak pressure $p_0$ under normal load $P$. Because the tangential load stays below the
sliding limit $\mu P$, the contact does not slide as a whole. Instead it splits into an
inner stick zone $r\le c$ and an outer slip annulus $c\le r\le a$. The stick radius $c$
shrinks as $Q$ approaches $\mu P$ \citep{mindlin1949compliance,johnson1985contact}:
% code: postprocessing/contact_fields.py::cattaneo_stick_radius
\begin{equation}\label{eq:cattaneo-c}
  \frac{c}{a} = \left(1-\frac{Q}{\mu P}\right)^{1/3}.
\end{equation}
The plane-strain solver of Section~\ref{sec:verif} verifies the two-dimensional line-contact
reduction, for which the stick radius follows the larger exponent
$c/a = (1-Q/\mu P)^{1/2}$; the axisymmetric $1/3$ form above is the spherical-contact reference.
The stick zone vanishes ($c\to0$) exactly when $Q$ reaches $\mu P$ and full sliding
begins. The tangential traction $q(r)$ superposes the full-sliding distribution on an
inner corrective term that cancels the slip displacement over the stick zone:
% code: postprocessing/contact_fields.py::cattaneo_traction
\begin{equation}\label{eq:cattaneo-q}
  q(r) =
  \begin{cases}
    \mu p_0\sqrt{1-r^2/a^2}, & c\le r\le a,\\[4pt]
    \mu p_0\!\left[\sqrt{1-r^2/a^2}-\dfrac{c}{a}\sqrt{1-r^2/c^2}\,\right], & r\le c.
  \end{cases}
\end{equation}
In the slip annulus the traction saturates at the Coulomb limit $\mu p(r)$, while inside
the stick zone the corrective term keeps it below that limit. Integrating $q$ over the
contact and relating it to the tangential motion gives the compliance $\delta_x$, the
tangential displacement under load:
% code: postprocessing/contact_fields.py::cattaneo_stick_radius
\begin{equation}\label{eq:cattaneo-delta}
  \delta_x = \frac{3\mu P}{16\,a\,G^\ast}\!\left[1-\left(1-\frac{Q}{\mu P}\right)^{2/3}\right],
  \qquad
  \frac{1}{G^\ast} = \frac{2-\nu_1}{4G_1}+\frac{2-\nu_2}{4G_2}.
\end{equation}
Here $G^\ast$ is the effective shear modulus, built from the shear moduli $G_1,G_2$ the
same way $E^\ast$ is built from the Young's moduli. The compliance is the frictional
counterpart of the normal force--approach law. It softens as $Q$ rises and the stick zone
gives way. We accept the numerical solution when $c/a$ agrees to $<1\%$, the integral
$\int q\,\mathrm{d}A = Q$ to $<1\%$, and $q(r)$ in \eqref{eq:cattaneo-q} agrees pointwise.

\subsection{CV-3 --- Brazilian disc}\label{app:brazilian}

A disc of diameter $D=2R$ and thickness $t$ is compressed diametrally by a force $P$. The
solution superposes two Flamant load-point charts at the poles $(0,\pm R)$ on a uniform
biaxial tension $T=2P/\pi Dt$ that cancels the spurious traction the point loads leave on
the rim \citep{timoshenko1970elasticity}. At the disc center the two charts combine into a
tensile horizontal stress and a compressive vertical stress:
% code: postprocessing/contact_fields.py::brazilian_field
\begin{equation}\label{eq:brazilian-center}
  \sigma_{xx} = +\frac{2P}{\pi D t} \;\;(\text{tensile}),
  \qquad
  \sigma_{yy} = -\frac{6P}{\pi D t} \;\;(\text{compressive}),
\end{equation}
so the horizontal tension at the center is one third the vertical compression, and this
tension splits the disc. The splitting tensile strength inferred from a failure load $P_f$
is therefore $\sigma_t = 2P_f/\pi D t$ (ASTM~D3967). Vertical equilibrium across the
diametral section gives an independent balance the field must satisfy:
% code: postprocessing/contact_fields.py::brazilian_field
\begin{equation}\label{eq:brazilian-balance}
  \int_{-R}^{R}\sigma_{yy}(x,0)\,\mathrm{d}x = -\frac{P}{t},
\end{equation}
that is, the vertical stress integrated across the diameter carries the full applied load
per unit thickness. The field must also satisfy the admissibility conditions
$\grad\!\cdot\bm{\sigma}=\bm{0}$ and $\grad^2(\mathrm{tr}\,\bm{\sigma})=0$ away from the
poles. Along the loaded diameter, a finite load patch replaces the point load, and the
distributed-pressure (Hondros) form \citep{hondros1959evaluation} gives, for a contact
half-angle $2\alpha$ subtended by the patch at the center,
% code: postprocessing/contact_fields.py::brazilian_field
\begin{equation}\label{eq:hondros}
  \sigma_{xx}(0,y) = \frac{2P}{\pi D t\,\alpha}\!\left[\,
    \frac{(1-r^2/R^2)\sin 2\alpha}{1-2r^2\cos2\alpha/R^2+r^4/R^4}
    -\arctan\!\Big(\frac{1+r^2/R^2}{1-r^2/R^2}\tan\alpha\Big)\right],
\end{equation}
with $r=|y|$ the distance from the center along the loaded diameter. As the patch shrinks
($\alpha\to0$), \eqref{eq:hondros} reduces to the center value \eqref{eq:brazilian-center}
and removes the point-load singularity. We require the center stresses
\eqref{eq:brazilian-center} to $<1\%$ and the global balance \eqref{eq:brazilian-balance}
to $<0.5\%$. The diametral compliance under a true point load diverges logarithmically. A
finite contact half-width $b$ keeps it finite:
% code: postprocessing/contact_fields.py::brazilian_contact_halfwidth
\begin{equation}\label{eq:brazilian-compliance}
  \Delta_D = \frac{4P}{\pi E t}\!\left[\ln\frac{2R}{b}-\frac{1-\nu}{2}\right],
\end{equation}
where $\Delta_D$ is the change in diameter along the load line. The prefactor and the
$-(1-\nu)/2$ term are robust, while the logarithmic argument carries the calibration of the
contact width $b$.

\subsection{CV-4 --- Nine-disc equibiaxial unit cell}\label{app:ninedisc}

A $3\times3$ square array of identical discs of radius $R$ is confined isotropically. No
diagonal contacts form, because the diagonal spacing $2\sqrt2\,R$ exceeds the contact
spacing $2R$. The $D_4$ symmetry of the lattice makes every contact carry the same normal
force $N$, while reflection parity keeps friction inactive. Each disc then sees four equal
inward loads $N$ at its north, south, east, and west poles, so the unit cell is the
superposition of two perpendicular Brazilian solutions, that is, four Flamant charts and
two uniform-tension charts. Adding the two perpendicular Brazilian fields makes the center
state an equibiaxial in-plane compression:
% code: postprocessing/contact_fields.py::nine_disc_unit_cell_field
\begin{equation}\label{eq:ninedisc-center}
  \sigma_{xx} = \sigma_{yy} = -\frac{2N}{\pi R t}, \qquad \sigma_{xy}=0,
\end{equation}
that is, equal compression in both directions and no shear. In the rigid-plate limit, the
numerical solution must reproduce this to $<1\%$ with an interior field identical across
the nine discs. The force--compression relation then closes the system through the
per-contact closing displacement $\Delta_x(N)$, the amount each contact shortens under
force $N$:
% code: postprocessing/contact_fields.py::nine_disc_unit_cell_field
\begin{equation}\label{eq:ninedisc-force}
  3\,\Delta_x(N) = 2d, \qquad
  \Delta_x(N) = \frac{4N}{\pi E t}\ln\frac{2R}{b}
  + \frac{N}{\pi E t}\big[(2-\pi)\nu + \pi - 6\big],
  \qquad
  b = 2\sqrt{\frac{N R_{\rm rel}}{\pi E^\ast}},
\end{equation}
where $d$ is the boundary compression and $b$ the Hertz half-width of a disc--disc
contact; the factor of three counts the three contacts a load crosses along one row. For
$R=1.7$~mm, $E=100$~GPa, $\nu=0.3$, and $t=1$~mm, the relation gives
$N(d{=}10^{-4}\,\text{mm}) \approx 0.887$~N/mm. Equation~\eqref{eq:ninedisc-center} holds
exactly only for rigid confining plates. Deformable plates widen and soften the disc--plate
contacts by $\sqrt2$, so the interior discs differ by a few percent near their boundary
contacts, and any initial gaps make the contacts engage progressively. We record these
caveats with each numerical comparison.

\begin{table}
\centering
\caption{Closed-form acceptance targets for the code-verification benchmarks CV-1--CV-4.
The numpy evaluators in \code{postprocessing/contact\_fields.py} are the single source of
truth, and the symbolic provenance is \code{docs/hertz\_derivation/}.}
\label{tab:closedforms}
\begin{tabular}{@{}lll@{}}
\toprule
Quantity & Closed form & Tolerance \\
\midrule
Hertz contact radius      & $a=(3FR/4E^\ast)^{1/3}$                 & $<1\%$ \\
Hertz peak pressure       & $p_0=3F/2\pi a^2$                       & $<1\%$ \\
Hertz approach            & $\delta=a^2/R$                          & $<1\%$ \\
Hertz force--approach     & $F=\tfrac43 E^\ast\sqrt{R}\,\delta^{3/2}$ & $<1\%$ \\
Pressure integral         & $\int p\,\mathrm{d}A=F$                 & $<0.5\%$ \\
Subsurface max shear      & $z/a=0.48$, $\tau_{\max}=0.31\,p_0$     & $\pm2\%$ \\
Line-contact half-width   & $a=2\sqrt{P'R/\pi E^\ast}$              & $<1\%$ \\
Cattaneo stick radius     & $c/a=(1-Q/\mu P)^{1/3}$                 & $<1\%$ \\
Brazilian center          & $\sigma_{xx}=+2P/\pi Dt,\ \sigma_{yy}=-6P/\pi Dt$ & $<1\%$ \\
Brazilian global balance  & $\int_{-R}^{R}\sigma_{yy}(x,0)\,\mathrm{d}x=-P/t$ & $<0.5\%$ \\
Nine-disc center          & $\sigma_{xx}=\sigma_{yy}=-2N/\pi Rt$    & $<1\%$ \\
\bottomrule
\end{tabular}
\end{table}
